\section{Discussion}
\label{sec:discussion}

\para{What do the results imply for a theory of simulation?} The evidence supports a middle position. \gptfourone is not just emitting one default answer. Persona conditioning improves \socialiqa accuracy and expands the set of ultimatum policies, which shows some responsiveness to latent factors. At the same time, the added diversity is narrow, structured, and biased toward compliance. This means the model behaves more like a controllable generator of behavioral manifolds than a simulator with full human-outcome coverage.

\para{Surface plausibility, fidelity, and coverage should be evaluated separately.} The experiments suggest a practical three-level view. Surface plausibility is easy: the model produces coherent human-like justifications in every condition. Local fidelity is harder but still attainable: on \socialiqa, persona prompting improves performance rather than harming it. Distributional coverage is the bottleneck. The model broadens support from 2 to 5 ultimatum profiles under persona prompting, but those policies remain clustered around very low acceptance thresholds. A theory of simulation that conflates these three levels will overstate current capabilities.

\para{Error analysis clarifies the failure modes.} In \socialiqa, the main errors are semantically narrow. The model often chooses a plausible affective state when the benchmark expects a more specific trait label, or it picks a broadly reasonable answer instead of the intended social inference. In the ultimatum task, the failures are structural. Demographic conditioning barely changes the policy family, and persona conditioning mainly introduces profiles that are even more accepting or only slightly stricter. This pattern matches prior concerns about average-person collapse and rationality bias \cite{liu2025rational,chen2026omnibehavior}.

\para{Limitations.} Our study tests only one model family and two short-horizon tasks. \socialiqa is likely affected by pretraining exposure, so it is not a pure simulation benchmark. The local OpinionQA mirror lacked the full human-response tables, which prevented a stronger subgroup-alignment analysis. The ultimatum design is lightweight and does not use richer memory or belief-desire-intention scaffolds of the kind explored in prior work \cite{xie2024trust,park2023generative}. Finally, the API behavior reflects \gptfourone as accessed on May 5, 2026 and may drift over time.

\para{Broader implications.} The main operational lesson is that LLMs should not be treated as drop-in substitutes for human populations in settings where subgroup structure or tail risk matters. They are more defensibly used as controllable generators that can help probe hypotheses, prototype interaction policies, or stress-test a bounded region of a behavior space. Stronger claims require separate validation for fidelity, subgroup structure, and tail coverage on the target domain.
